\section{Reconstruction with unknown smooth remainder functions}
\label{sec:v8-nuisance}
The exact inverse in Section~\ref{sec:v8-fixed} uses the known nonlinear
family at positive offset. To separate that assumption from the
acquisition of leading invariants, we now enlarge the observation
model explicitly. The underlying parameter is still $(g,e)\in K$ and
its leading amplitudes are still the physical $C_j(g,e)$. The nuisance
functions below need not come from that exact analytic support family.

Fix $m\ge1$, constants $0<c_*<C_*<\infty$, and a sufficiently small
$d_*>0$. For each $j=1,2,3$, let the unknown function $H_j$ satisfy
\begin{equation}\label{eq:v8-envelope}
 \begin{gathered}
 H_j\in C^m([0,d_*]),\quad c_*\le H_j\le C_*,\quad
       \|H_j\|_{C^m}\le C_*,\quad H_j(0)=C_j(g,e),\\
 p_j(t)=(t-jg)_+^2H_j((t-jg)_+),\qquad t-jg\le d_*.
 \end{gathered}
\end{equation}
Choose $C_*d_*^2<1$. Each query returns an independent Bernoulli
variable with this probability. The fixed neighborhood $K$, the common
bounds, and a fixed interval $[L_0,U_0]$ containing the possible gaps
are known. No value or derivative of $H_j$ away from zero is supplied.
The class can be chosen to contain the exact physical finite-offset
probabilities on a smaller compact neighborhood. It is an observation
envelope: we do not assert that every arbitrary nuisance triple in it
is realized by a single billiard table. The target area and curvatures
are defined from $(g,e)$ by the same physical leading inverse.

\begin{theorem}[Uniform count recovery over a smooth envelope]
\label{thm:v8-nuisance}
Let $0<U_0-L_0=W<d_*/12$. For fixed $m$ and the preceding bounds,
there are $h_m,C_m>0$ such that for $0<h<h_m$ and $0<\eta<1/4$
a Borel count-only procedure satisfies, uniformly over
\eqref{eq:v8-envelope},
\begin{equation}\label{eq:v8-nuisance-errors}
 \Pr\left\{
 \begin{array}{l}
 |\widehat g-g|\le C_mh^{m+1},\\
 \|\widehat e-e\|_\infty+|\widehat A-A|\le C_mh^m,\\
 \operatorname{dist}_{\rm match}(\widehat\kappa,\kappa)
                                      \le C_mh^{m/3}
 \end{array}\right\}\ge1-\eta.
\end{equation}
Its number of preparations is bounded on every history by
\begin{equation}\label{eq:v8-nuisance-cost}
                       C_mh^{-(2m+2)}\log(C_m/\eta).
\end{equation}
It uses only flight numbers $1,2,3$ and programmed times. It supplies
its own accuracy-dependent bracket from the fixed interval, and
requires no exact finite-offset formula, contact positions, labels,
or independently measured area.
\end{theorem}
\begin{proof}
We give a deterministic-budget form of the calibration and extrapolation
construction to specify its dependence on the nuisance functions.

First use the bracket search of Lemma~\ref{lem:v7-bracket} with
failure allowance $\eta/3$. Its proof needs only zero probability
below onset and the lower bound $p_1(g+d)\ge c_*d^2$, both part of
\eqref{eq:v8-envelope}. It returns a midpoint $g_0\in[L_0,U_0]$ and,
on its good event, $|g-g_0|\le h/4$. Its all-history cost is
$C h^{-2}\log((2+\log(W/h))/\eta)$.

Next take a predetermined
\[
                  n=\left\lceil C_mh^{-(2m+2)}
                                      \log(C_m/\eta)\right\rceil
\]
fresh preparations at each pilot time $g_0+lh$,
$1\le l\le m+1$. On bracket success all actual offsets lie between
$3h/4$ and $(m+5/4)h$, and the probabilities are between positive
multiples of $h^2$. The Bernoulli concentration bound proved in
Lemma~\ref{lem:v6-amplitudes}, divided by those probabilities, gives
relative errors at most $c_mh^m$ simultaneously, with conditional
probability at least $1-\eta/3$, after increasing $C_m$.

Interpolate the square roots of the empirical fractions by a polynomial
of degree at most $m$ in the variable $z=t-g_0$ at nodes $lh$.
Use its unique-zero-or-zero convention on $[-h/2,h/2]$ and call the
result $\widehat\tau$. This rule is defined and Borel also when a
fraction is zero. Put $\widehat g=g_0+\widehat\tau$. On every history
$|\widehat g-g_0|\le h/2$. On the good pilot event, write
\[
 d\sqrt{H_1(d)}=\sqrt{C_1}\,d
                      \sqrt{H_1(d)/C_1}.
\]
The second square-root factor has value one at zero and bounded
$C^m$ norm, uniformly over the envelope. Multiplying its degree-$(m-1)$
Taylor polynomial by $d$ gives a polynomial of degree $m$ with a
simple zero at onset. The fixed-node Lagrange bounds and the root
argument of Lemma~\ref{lem:v5-calibration} therefore give
$|\widehat g-g|\le C_mh^{m+1}$. Only right derivatives through order
$m$ were used; no probability was evaluated below onset for this
Taylor argument.

For the last stage put $L=4$ and take $n$ fresh preparations at
$j\widehat g+Llh$ for $j=1,2,3$ and $l=1,\ldots,m$. Use
\[
 \omega_l=(-1)^{l-1}\binom ml,\qquad
 \widehat C_j=\sum_{l=1}^m\omega_l
                              \frac{\widehat p_{j,l}}{(Llh)^2}.
\]
On bracket success every actual offset of this stage lies in a fixed
positive interval of multiples of $h$, even if the pilot estimates
are inaccurate. On the two good events, its displacement from $Llh$
is at most $3C_mh^{m+1}$. The derivative of $d^2H_j(d)$ is $O_m(h)$
on these intervals, so the resulting normalized timing error is
$O_m(h^m)$. The identities
\[
          \sum_l\omega_l=1,\qquad
          \sum_l\omega_l l^r=0\quad(1\le r<m)
\]
and the bounded $m$th derivative of $H_j$ give extrapolation bias
$O_m(h^m)$. Conditional Bernoulli concentration on these fresh
preparations gives the same stochastic order, with additional failure
probability at most $\eta/3$. Hence
$\max_j|\widehat C_j-C_j|\le C_mh^m$ on their intersection.

Fit $(\widehat g,\widehat C_1,\widehat C_2,\widehat C_3)$ on the
fixed compact physical leading-data set using
Lemma~\ref{lem:v7-selection}. The jointly Lipschitz coefficient inverse
of Theorem~\ref{thm:v6-three} gives coefficient and area errors
$O_m(h^m)$, and its cubic root bound gives curvature error $O_m(h^{m/3})$.
Retain the pilot $\widehat g$ as the gap estimate; the fitted physical
parameter and the pilot need not have identical gaps. This distinction
preserves the finer asserted timing error.

All queries are legitimate even after a failed search. The returned
$g_0$ remains in $[L_0,U_0]$ and
$|\widehat g-g_0|\le h/2$. Thus pilot positive offsets are at most
$W+(m+1)h$, and final positive offsets at most
$3(W+h/2)+4mh$. The fixed restriction on $W$ and small enough $h_m$
place both in $[0,d_*]$; negative offsets are deterministic failures.
No waiting time depends on seeing a success. The pilot and final stage
have exactly $(m+1+3m)n$ preparations on every history. As in
Theorem~\ref{thm:v7-fixed-bracket}, their bound absorbs the bracket
search cost because $h^{2m}\log(2+\log(W/h))$ is bounded. The three
conditional failure allowances sum to at most $\eta$, proving the
result.
\end{proof}

Taking $h\asymp\varepsilon^{3/m}$ yields curvature error
$O(\varepsilon)$ and gap error $O(\varepsilon^{3+3/m})$ with the
preparation power $6+6/m$. This result is uniform over unknown smooth
higher-order remainders. The theorem gives an upper bound; the matching envelope lower bound,
under the additional strict interior condition, is proved in
Theorem~\ref{thm:v9-envelope-minimax}. In the exact analytic
subexperiment, the same joint accuracy has the lower bound of
Theorem~\ref{thm:v8-joint-minimax}; for curvature loss alone its
fixed-window benchmark has power six instead. The two experiments
therefore differ in the information supplied to the estimator, even
when their stated joint guarantees have the same preparation order.
